% Section 7: Metadata and global index file fields

\section{Metadata and global index file fields}
\label{sec:metadata-and-global-index-file-fields}


\subsection{Reprojected image and mosaic metadata params}
\label{sec:reprojected-image-and-mosaic-metadata-params}

Units are uniform throughout this section and throughout the bundle. All angles, longitudes, and true anomalies are in degrees; all radii and radial distances are in km; radial resolution is in km/pixel and longitudinal resolution in degrees/pixel; coverage is given in percent; and all times given as \code{tdb} (including \code{rings:\allowbreak{}observed\_\allowbreak{}event\_\allowbreak{}tdb}) are ephemeris times in seconds past the J2000 epoch, 2000-01-01T12:00:00 Barycentric Dynamical Time (TDB) — the same quantity SPICE calls \textit{ET}. Times given as \code{date\_\allowbreak{}time} are Coordinated Universal Time (UTC).

The following fields are tabulated for each valid corotating longitude in each reprojected image (\code{IMGID\_\allowbreak{}reproj\_\allowbreak{}img\_\allowbreak{}metadata\_\allowbreak{}params.\allowbreak{}tab}), mosaic (\code{OBSID\_\allowbreak{}mosaic\_\allowbreak{}metadata\_\allowbreak{}params.\allowbreak{}tab}), and background-subtracted mosaic (\code{OBSID\_\allowbreak{}mosaic\_\allowbreak{}bkg\_\allowbreak{}sub\_\allowbreak{}metadata\_\allowbreak{}params.\allowbreak{}tab}).

\begingroup\small
\begin{longtable}{|>{\raggedright\arraybackslash}p{0.255\linewidth}|>{\raggedright\arraybackslash}p{0.665\linewidth}|}
\hline
\code{rings:\allowbreak{}corotating\_\allowbreak{}ring\_\allowbreak{}longitude} & The corotating longitude for which values are being provided \\
\hline
\code{image\_\allowbreak{}index} & (Mosaics only) The number of the reprojected image (from \code{OBSID\_\allowbreak{}mosaic\_\allowbreak{}metadata\_\allowbreak{}src\_\allowbreak{}imgs.\allowbreak{}tab} or \code{OBSID\_\allowbreak{}mosaic\_\allowbreak{}bkg\_\allowbreak{}sub\_\allowbreak{}metadata\_\allowbreak{}src\_\allowbreak{}imgs.\allowbreak{}tab}) that provided data for this corotating longitude \\
\hline
\code{rings:\allowbreak{}observed\_\allowbreak{}event\_\allowbreak{}tdb} & The mid-time of the observation that provided data for this corotating longitude (constant for one image) \\
\hline
\code{rings:\allowbreak{}inertial\_\allowbreak{}ring\_\allowbreak{}longitude} & The inertial longitude of the part of the source image that provided data for this corotating longitude \\
\hline
\code{rings:\allowbreak{}radial\_\allowbreak{}resolution} & The radial resolution of the part of the source image that provided data for this corotating longitude \\
\hline
\code{rings:\allowbreak{}longitudinal\_\allowbreak{}resolution} & The angular longitudinal resolution of the part of the source image that provided data for this corotating longitude \\
\hline
\code{rings:\allowbreak{}incidence\_\allowbreak{}angle} & The incidence angle of the ring plane at the time of the source observation (constant for one image) \\
\hline
\code{rings:\allowbreak{}phase\_\allowbreak{}angle} & The phase angle of the part of the source image that provided data for this corotating longitude \\
\hline
\code{rings:\allowbreak{}emission\_\allowbreak{}angle} & The emission angle of the part of the source image that provided data for this corotating longitude \\
\hline
\code{core\_\allowbreak{}radius} & The radius of the ideal F ring core from the orbit model at this corotating longitude \\
\hline
\code{longitude\_\allowbreak{}ascending\_\allowbreak{}node} & The longitude of the ascending node (relative to J2000) from the orbit model at the time of the source observation (constant for one image) \\
\hline
\code{longitude\_\allowbreak{}pericenter} & The longitude of pericenter (relative to J2000) from the orbit model at the time of the source observation (constant for one image) \\
\hline
\code{true\_\allowbreak{}anomaly} & The true anomaly from the orbit model corresponding to this corotating longitude \\
\hline
\code{corotating\_\allowbreak{}longitude\_\allowbreak{}prometheus} & The corotating longitude of Prometheus at the time of the source image, in the same reference frame used for the F ring core (constant for one image) \\
\hline
\code{radius\_\allowbreak{}prometheus} & The orbital radius (distance from Saturn center) of Prometheus at the time of the original observation \\
\hline
\code{corotating\_\allowbreak{}longitude\_\allowbreak{}pandora} & The corotating longitude of Pandora at the time of the source image, in the same reference frame used for the F ring core (constant for one image) \\
\hline
\code{radius\_\allowbreak{}pandora} & The orbital radius (distance from Saturn center) of Pandora at the time of the original observation \\
\hline
\end{longtable}
\endgroup


\subsection{Reprojected image global index}
\label{sec:reprojected-image-global-index}

The following fields appear in \code{global\_\allowbreak{}reproj\_\allowbreak{}img\_\allowbreak{}index.\allowbreak{}tab}.

\begingroup\small
\begin{longtable}{|>{\raggedright\arraybackslash}p{0.313\linewidth}|>{\raggedright\arraybackslash}p{0.607\linewidth}|}
\hline
\code{pds:\allowbreak{}logical\_\allowbreak{}identifier} & The PDS4 logical identifier, \textit{e.g.} \code{urn:\allowbreak{}nasa:\allowbreak{}pds:\allowbreak{}cassini\_\allowbreak{}iss\_\allowbreak{}fring\_\allowbreak{}mosaics\_\allowbreak{}rsfrench2025:\allowbreak{}}\newline \code{data\_\allowbreak{}reproj\_\allowbreak{}img:\allowbreak{}1874525875w\_\allowbreak{}reproj\_\allowbreak{}img} \\
\hline
\code{cassini:\allowbreak{}observation\_\allowbreak{}id} & The Cassini observation name of the source image, \textit{e.g.} \code{IOSIC\_\allowbreak{}276RB\_\allowbreak{}COMPLITB3001\_\allowbreak{}SI} \\
\hline
\code{file\_\allowbreak{}spec} & The path to the reprojected image label file relative to the bundle root, \textit{e.g.} \code{data\_\allowbreak{}reproj\_\allowbreak{}img/\allowbreak{}iosic\_\allowbreak{}276rb\_\allowbreak{}complitb3001\_\allowbreak{}si/\allowbreak{}}\newline \code{1874525875w\_\allowbreak{}reproj\_\allowbreak{}img.\allowbreak{}lblx} \\
\hline
\code{pds:\allowbreak{}start\_\allowbreak{}date\_\allowbreak{}time} & The start time of the source image exposure, \textit{e.g.} \code{2017-05-26T20:\allowbreak{}29:\allowbreak{}10Z} \\
\hline
\code{pds:\allowbreak{}stop\_\allowbreak{}date\_\allowbreak{}time} & The stop time of the source image exposure \\
\hline
\code{cassini:\allowbreak{}spacecraft\_\allowbreak{}clock\_\allowbreak{}start\_\allowbreak{}count} & The starting spacecraft clock count of the source image without the partition number (which is always \code{1/\allowbreak{}}) \\
\hline
\code{cassini:\allowbreak{}spacecraft\_\allowbreak{}clock\_\allowbreak{}stop\_\allowbreak{}count} & The ending spacecraft clock count of the source image without the partition number (which is always \code{1/\allowbreak{}}) \\
\hline
\code{num\_\allowbreak{}valid\_\allowbreak{}longitudes} & The number of discrete longitudes (radial slices) that contain valid data; each radial slice represents 0.02° of corotating longitude \\
\hline
\code{percent\_\allowbreak{}coverage} & The percentage of 360° containing valid data; computed as \code{num\_\allowbreak{}valid\_\allowbreak{}longitudes} / 18,000 × 100 \\
\hline
\code{rings:\allowbreak{}minimum\_\allowbreak{}corotating\_\allowbreak{}ring\_\allowbreak{}longitude} & The minimum corotating longitude, representing the left edge of the reprojected image \\
\hline
\code{rings:\allowbreak{}maximum\_\allowbreak{}corotating\_\allowbreak{}ring\_\allowbreak{}longitude} & The maximum corotating longitude, representing the right edge of the reprojected image; note that it is possible for min > max, in which case the corotating range wraps around at 360° \\
\hline
\code{rings:\allowbreak{}minimum\_\allowbreak{}inertial\_\allowbreak{}ring\_\allowbreak{}longitude} & The minimum inertial longitude from the source image used to create the reprojected image \\
\hline
\code{rings:\allowbreak{}maximum\_\allowbreak{}inertial\_\allowbreak{}ring\_\allowbreak{}longitude} & The maximum inertial longitude from the source image used to create the reprojected image; because this is a single exposure, the corotating and inertial longitude ranges have the same width \\
\hline
\code{rings:\allowbreak{}mean\_\allowbreak{}phase\_\allowbreak{}angle} & The mean (across all valid longitudes) phase angle \\
\hline
\code{rings:\allowbreak{}minimum\_\allowbreak{}phase\_\allowbreak{}angle} & The minimum (across all valid longitudes) phase angle \\
\hline
\code{rings:\allowbreak{}maximum\_\allowbreak{}phase\_\allowbreak{}angle} & The maximum (across all valid longitudes) phase angle \\
\hline
\code{rings:\allowbreak{}mean\_\allowbreak{}incidence\_\allowbreak{}angle} & The ring plane incidence angle at the time of the observation. The incidence angle does not vary with longitude and changes only slowly with time, so a single value is computed for the whole product; \code{rings:\allowbreak{}minimum\_\allowbreak{}incidence\_\allowbreak{}angle} and \code{rings:\allowbreak{}maximum\_\allowbreak{}incidence\_\allowbreak{}angle} are present in the label and carry that same value \\
\hline
\code{rings:\allowbreak{}mean\_\allowbreak{}emission\_\allowbreak{}angle} & The mean (across all valid longitudes) emission angle; the angle is < 90° on the lit side of the rings \\
\hline
\code{rings:\allowbreak{}minimum\_\allowbreak{}emission\_\allowbreak{}angle} & The minimum (across all valid longitudes) emission angle \\
\hline
\code{rings:\allowbreak{}maximum\_\allowbreak{}emission\_\allowbreak{}angle} & The maximum (across all valid longitudes) emission angle \\
\hline
\code{rings:\allowbreak{}mean\_\allowbreak{}radial\_\allowbreak{}resolution} & The mean (across all valid longitudes) radial resolution in km/pixel; radial resolution is the resolution in the radial direction of the source image used to populate each reprojected image pixel \\
\hline
\code{rings:\allowbreak{}minimum\_\allowbreak{}radial\_\allowbreak{}resolution} & The minimum (across all valid longitudes) radial resolution in km/pixel \\
\hline
\code{rings:\allowbreak{}maximum\_\allowbreak{}radial\_\allowbreak{}resolution} & The maximum (across all valid longitudes) radial resolution in km/pixel \\
\hline
\code{rings:\allowbreak{}mean\_\allowbreak{}longitudinal\_\allowbreak{}resolution} & The mean (across all valid longitudes) angular longitudinal resolution in °/pixel; angular longitudinal resolution is the resolution in the longitude direction (expressed as an angle) of the source image used to populate each reprojected image pixel \\
\hline
\code{rings:\allowbreak{}minimum\_\allowbreak{}longitudinal\_\allowbreak{}resolution} & The minimum (across all valid longitudes) angular longitudinal resolution in °/pixel \\
\hline
\code{rings:\allowbreak{}maximum\_\allowbreak{}longitudinal\_\allowbreak{}resolution} & The maximum (across all valid longitudes) angular longitudinal resolution in °/pixel \\
\hline
\code{rings:\allowbreak{}minimum\_\allowbreak{}ring\_\allowbreak{}radius} & The minimum (across all valid longitudes) ring radius (distance from Saturn center) in the original image used to provide data at any valid longitude \\
\hline
\code{rings:\allowbreak{}maximum\_\allowbreak{}ring\_\allowbreak{}radius} & The maximum (across all valid longitudes) ring radius (distance from Saturn center) in the original image used to provide data at any valid longitude \\
\hline
\code{mean\_\allowbreak{}core\_\allowbreak{}radius} & The mean (across all valid longitudes) ring radius (distance from Saturn center) of the ideal F ring core from the orbit model (Section~\ref{sec:f-ring-orbit}) \\
\hline
\code{minimum\_\allowbreak{}core\_\allowbreak{}radius} & The minimum (across all valid longitudes) ring radius (distance from Saturn center) of the F ring core from the orbit model \\
\hline
\code{maximum\_\allowbreak{}core\_\allowbreak{}radius} & The maximum (across all valid longitudes) ring radius (distance from Saturn center) of the F ring core from the orbit model \\
\hline
\code{longitude\_\allowbreak{}ascending\_\allowbreak{}node} & The longitude of the ascending node of the F ring (relative to J2000) from the orbit model at the time of the original observation (Section~\ref{sec:f-ring-orbit}) \\
\hline
\code{longitude\_\allowbreak{}pericenter} & The longitude of pericenter of the F ring (relative to J2000) from the orbit model at the time of the original observation (Section~\ref{sec:f-ring-orbit}) \\
\hline
\code{minimum\_\allowbreak{}true\_\allowbreak{}anomaly} & The minimum (across all valid longitudes) true anomaly from the orbit model (Section~\ref{sec:f-ring-orbit}) \\
\hline
\code{maximum\_\allowbreak{}true\_\allowbreak{}anomaly} & The maximum (across all valid longitudes) true anomaly from the orbit model (Section~\ref{sec:f-ring-orbit}) \\
\hline
\code{corotating\_\allowbreak{}longitude\_\allowbreak{}prometheus} & The corotating longitude of Prometheus at the time of the original observation in the same reference frame used for the F ring core \\
\hline
\code{radius\_\allowbreak{}prometheus} & The orbital radius (distance from Saturn center) of Prometheus at the time of the original observation \\
\hline
\code{corotating\_\allowbreak{}longitude\_\allowbreak{}pandora} & The corotating longitude of Pandora at the time of the original observation in the same reference frame used for the F ring core \\
\hline
\code{radius\_\allowbreak{}pandora} & The orbital radius (distance from Saturn center) of Pandora at the time of the original observation \\
\hline
\code{pds:\allowbreak{}creation\_\allowbreak{}date\_\allowbreak{}time} & The date and time the reprojected image product was created \\
\hline
\code{nav\_\allowbreak{}quality} & The navigation grade of the mosaic belonging to this image’s observation. It is the overall grade for that group of images, not an assessment of this one image, and every reprojected image of the same observation carries the same value — including those that did not contribute to the mosaic (Section~\ref{sec:reprojected-images-not-used-in-a-mosaic}). Values are G = Good, F = Fair, P = Poor (Section~\ref{sec:navigation}) \\
\hline
\code{notes} & Notes about the mosaic belonging to this image’s observation (Section~\ref{sec:the-global-index-files}) \\
\hline
\end{longtable}
\endgroup


\subsection{Mosaic and background-subtracted mosaic global indexes}
\label{sec:mosaic-global-index}

The following fields appear in \code{global\_\allowbreak{}mosaic\_\allowbreak{}index.\allowbreak{}tab} and \code{global\_\allowbreak{}mosaic\_\allowbreak{}bkg\_\allowbreak{}sub\_\allowbreak{}index.\allowbreak{}tab}. Rows marked “background-subtracted mosaic only” are present only in the latter.

\begingroup\small
\begin{longtable}{|>{\raggedright\arraybackslash}p{0.313\linewidth}|>{\raggedright\arraybackslash}p{0.607\linewidth}|}
\hline
\code{pds:\allowbreak{}logical\_\allowbreak{}identifier} & The PDS4 logical identifier, \textit{e.g.} \code{urn:\allowbreak{}nasa:\allowbreak{}pds:\allowbreak{}cassini\_\allowbreak{}iss\_\allowbreak{}fring\_\allowbreak{}mosaics\_\allowbreak{}rsfrench2025:\allowbreak{}data\_\allowbreak{}mosaic:\allowbreak{}iosic\_\allowbreak{}276rb\_\allowbreak{}complitb3001\_\allowbreak{}si\_\allowbreak{}mosaic} \\
\hline
\code{cassini:\allowbreak{}observation\_\allowbreak{}id} & The Cassini observation name of the source images, \textit{e.g.} \code{IOSIC\_\allowbreak{}276RB\_\allowbreak{}COMPLITB3001\_\allowbreak{}SI} \\
\hline
\code{file\_\allowbreak{}spec} & The path to the mosaic label file relative to the bundle root, \textit{e.g.} \code{data\_\allowbreak{}mosaic/\allowbreak{}iosic\_\allowbreak{}276rb\_\allowbreak{}complitb3001\_\allowbreak{}si/\allowbreak{}iosic\_\allowbreak{}276rb\_\allowbreak{}complitb3001\_\allowbreak{}si\_\allowbreak{}mosaic.\allowbreak{}lblx} \\
\hline
\code{pds:\allowbreak{}start\_\allowbreak{}date\_\allowbreak{}time} & The start time of the earliest source image exposure, \textit{e.g.} \code{2017-05-26T20:\allowbreak{}29:\allowbreak{}10Z} \\
\hline
\code{pds:\allowbreak{}stop\_\allowbreak{}date\_\allowbreak{}time} & The stop time of the latest source image exposure \\
\hline
\code{cassini:\allowbreak{}spacecraft\_\allowbreak{}clock\_\allowbreak{}start\_\allowbreak{}count} & The starting spacecraft clock count of the earliest source image without the partition number (which is always \code{1/\allowbreak{}}) \\
\hline
\code{cassini:\allowbreak{}spacecraft\_\allowbreak{}clock\_\allowbreak{}stop\_\allowbreak{}count} & The ending spacecraft clock count of the latest source image without the partition number (which is always \code{1/\allowbreak{}}) \\
\hline
\code{num\_\allowbreak{}valid\_\allowbreak{}longitudes} & The number of discrete longitudes (radial slices) that contain valid data; each radial slice represents 0.02° of corotating longitude \\
\hline
\code{percent\_\allowbreak{}coverage} & The percentage of 360° containing valid data; computed as \code{num\_\allowbreak{}valid\_\allowbreak{}longitudes} / 18,000 × 100 \\
\hline
\code{rings:\allowbreak{}minimum\_\allowbreak{}corotating\_\allowbreak{}ring\_\allowbreak{}longitude} & The minimum corotating longitude, representing the left edge of the mosaic \\
\hline
\code{rings:\allowbreak{}maximum\_\allowbreak{}corotating\_\allowbreak{}ring\_\allowbreak{}longitude} & The maximum corotating longitude, representing the right edge of the mosaic; note that it is possible for min > max, in which case the corotating range wraps around at 360° \\
\hline
\code{rings:\allowbreak{}minimum\_\allowbreak{}inertial\_\allowbreak{}ring\_\allowbreak{}longitude} & The minimum inertial longitude from the source images used to create the mosaic \\
\hline
\code{rings:\allowbreak{}maximum\_\allowbreak{}inertial\_\allowbreak{}ring\_\allowbreak{}longitude} & The maximum inertial longitude from the source images used to create the mosaic \\
\hline
\code{rings:\allowbreak{}mean\_\allowbreak{}phase\_\allowbreak{}angle} & The mean (across all valid longitudes) phase angle \\
\hline
\code{rings:\allowbreak{}minimum\_\allowbreak{}phase\_\allowbreak{}angle} & The minimum (across all valid longitudes) phase angle \\
\hline
\code{rings:\allowbreak{}maximum\_\allowbreak{}phase\_\allowbreak{}angle} & The maximum (across all valid longitudes) phase angle \\
\hline
\code{rings:\allowbreak{}mean\_\allowbreak{}incidence\_\allowbreak{}angle} & The ring plane incidence angle during the times the source images were taken. The incidence angle does not vary with longitude and changes only slowly with time, so a single value is computed for the whole mosaic; \code{rings:\allowbreak{}minimum\_\allowbreak{}incidence\_\allowbreak{}angle} and \code{rings:\allowbreak{}maximum\_\allowbreak{}incidence\_\allowbreak{}angle} are present in the label and carry that same value \\
\hline
\code{rings:\allowbreak{}mean\_\allowbreak{}emission\_\allowbreak{}angle} & The mean (across all valid longitudes) emission angle; the angle is < 90° on the lit side of the rings \\
\hline
\code{rings:\allowbreak{}minimum\_\allowbreak{}emission\_\allowbreak{}angle} & The minimum (across all valid longitudes) emission angle \\
\hline
\code{rings:\allowbreak{}maximum\_\allowbreak{}emission\_\allowbreak{}angle} & The maximum (across all valid longitudes) emission angle \\
\hline
\code{rings:\allowbreak{}mean\_\allowbreak{}radial\_\allowbreak{}resolution} & The mean (across all valid longitudes) radial resolution in km/pixel; radial resolution is the resolution in the radial direction of the source image used to populate each mosaic pixel \\
\hline
\code{rings:\allowbreak{}minimum\_\allowbreak{}radial\_\allowbreak{}resolution} & The minimum (across all valid longitudes) radial resolution in km/pixel \\
\hline
\code{rings:\allowbreak{}maximum\_\allowbreak{}radial\_\allowbreak{}resolution} & The maximum (across all valid longitudes) radial resolution in km/pixel \\
\hline
\code{rings:\allowbreak{}mean\_\allowbreak{}longitudinal\_\allowbreak{}resolution} & The mean (across all valid longitudes) angular longitudinal resolution in °/pixel; angular longitudinal resolution is the resolution in the longitude direction (expressed as an angle) of the source image used to populate each mosaic pixel \\
\hline
\code{rings:\allowbreak{}minimum\_\allowbreak{}longitudinal\_\allowbreak{}resolution} & The minimum (across all valid longitudes) angular longitudinal resolution in °/pixel \\
\hline
\code{rings:\allowbreak{}maximum\_\allowbreak{}longitudinal\_\allowbreak{}resolution} & The maximum (across all valid longitudes) angular longitudinal resolution in °/pixel \\
\hline
\code{rings:\allowbreak{}minimum\_\allowbreak{}ring\_\allowbreak{}radius} & The minimum (across all valid longitudes) ring radius (distance from Saturn center) in the original images used to provide data at any valid longitude \\
\hline
\code{rings:\allowbreak{}maximum\_\allowbreak{}ring\_\allowbreak{}radius} & The maximum (across all valid longitudes) ring radius (distance from Saturn center) in the original images used to provide data at any valid longitude \\
\hline
\code{mean\_\allowbreak{}core\_\allowbreak{}radius} & The mean (across all valid longitudes) ring radius (distance from Saturn center) of the ideal F ring core (Section~\ref{sec:f-ring-orbit}) \\
\hline
\code{minimum\_\allowbreak{}core\_\allowbreak{}radius} & The minimum (across all valid longitudes) ring radius (distance from Saturn center) of the F ring core \\
\hline
\code{maximum\_\allowbreak{}core\_\allowbreak{}radius} & The maximum (across all valid longitudes) ring radius (distance from Saturn center) of the F ring core \\
\hline
\code{mean\_\allowbreak{}longitude\_\allowbreak{}ascending\_\allowbreak{}node} & The mean longitude of the ascending node of the F ring (relative to J2000) from the orbit model across all of the source images (Section~\ref{sec:f-ring-orbit}) \\
\hline
\code{minimum\_\allowbreak{}longitude\_\allowbreak{}ascending\_\allowbreak{}node} & The minimum longitude of the ascending node of the F ring (relative to J2000) from the orbit model across all of the source images (Section~\ref{sec:f-ring-orbit}) \\
\hline
\code{maximum\_\allowbreak{}longitude\_\allowbreak{}ascending\_\allowbreak{}node} & The maximum longitude of the ascending node of the F ring (relative to J2000) from the orbit model across all of the source images (Section~\ref{sec:f-ring-orbit}) \\
\hline
\code{mean\_\allowbreak{}longitude\_\allowbreak{}pericenter} & The mean longitude of pericenter of the F ring (relative to J2000) from the orbit model across all of the source images (Section~\ref{sec:f-ring-orbit}) \\
\hline
\code{minimum\_\allowbreak{}longitude\_\allowbreak{}pericenter} & The minimum longitude of pericenter of the F ring (relative to J2000) from the orbit model across all of the source images (Section~\ref{sec:f-ring-orbit}) \\
\hline
\code{maximum\_\allowbreak{}longitude\_\allowbreak{}pericenter} & The maximum longitude of pericenter of the F ring (relative to J2000) from the orbit model across all of the source images (Section~\ref{sec:f-ring-orbit}) \\
\hline
\code{minimum\_\allowbreak{}true\_\allowbreak{}anomaly} & The minimum (across all valid longitudes) true anomaly from the orbit model (Section~\ref{sec:f-ring-orbit}) \\
\hline
\code{maximum\_\allowbreak{}true\_\allowbreak{}anomaly} & The maximum (across all valid longitudes) true anomaly from the orbit model (Section~\ref{sec:f-ring-orbit}) \\
\hline
\code{mean\_\allowbreak{}corotating\_\allowbreak{}longitude\_\allowbreak{}prometheus} & The mean corotating longitude of Prometheus in the same reference frame used for the F ring core across all source images \\
\hline
\code{minimum\_\allowbreak{}corotating\_\allowbreak{}longitude\_\allowbreak{}prometheus} & The minimum corotating longitude of Prometheus in the same reference frame used for the F ring core across all source images \\
\hline
\code{maximum\_\allowbreak{}corotating\_\allowbreak{}longitude\_\allowbreak{}prometheus} & The maximum corotating longitude of Prometheus in the same reference frame used for the F ring core across all source images \\
\hline
\code{mean\_\allowbreak{}radius\_\allowbreak{}prometheus} & The mean orbital radius (distance from Saturn center) of Prometheus across all source images \\
\hline
\code{minimum\_\allowbreak{}radius\_\allowbreak{}prometheus} & The minimum orbital radius (distance from Saturn center) of Prometheus across all source images \\
\hline
\code{maximum\_\allowbreak{}radius\_\allowbreak{}prometheus} & The maximum orbital radius (distance from Saturn center) of Prometheus across all source images \\
\hline
\code{mean\_\allowbreak{}corotating\_\allowbreak{}longitude\_\allowbreak{}pandora} & The mean corotating longitude of Pandora in the same reference frame used for the F ring core across all source images \\
\hline
\code{minimum\_\allowbreak{}corotating\_\allowbreak{}longitude\_\allowbreak{}pandora} & The minimum corotating longitude of Pandora in the same reference frame used for the F ring core across all source images \\
\hline
\code{maximum\_\allowbreak{}corotating\_\allowbreak{}longitude\_\allowbreak{}pandora} & The maximum corotating longitude of Pandora in the same reference frame used for the F ring core across all source images \\
\hline
\code{mean\_\allowbreak{}radius\_\allowbreak{}pandora} & The mean orbital radius (distance from Saturn center) of Pandora across all source images \\
\hline
\code{minimum\_\allowbreak{}radius\_\allowbreak{}pandora} & The minimum orbital radius (distance from Saturn center) of Pandora across all source images \\
\hline
\code{maximum\_\allowbreak{}radius\_\allowbreak{}pandora} & The maximum orbital radius (distance from Saturn center) of Pandora across all source images \\
\hline
\code{pds:\allowbreak{}creation\_\allowbreak{}date\_\allowbreak{}time} & The date and time the mosaic product was created \\
\hline
\code{nav\_\allowbreak{}quality} & The overall quality of the navigation across all source images; values are G = Good, F = Fair, P = Poor (Section~\ref{sec:navigation}) \\
\hline
\code{bkgnd\_\allowbreak{}quality} & (Background-subtracted mosaic only) The overall quality of the background model used for this mosaic; values are G = Good, F = Fair, P = Poor (Section~\ref{sec:background-modeling}) \\
\hline
\code{notes} & Notes about the mosaic (Section~\ref{sec:the-global-index-files}) \\
\hline
\code{num\_\allowbreak{}images} & The number of reprojected images used to create this mosaic \\
\hline
\code{min\_\allowbreak{}image\_\allowbreak{}name} & The name of the earliest source image used in this mosaic \\
\hline
\code{max\_\allowbreak{}image\_\allowbreak{}name} & The name of the latest source image used in this mosaic \\
\hline
\code{bkgnd\_\allowbreak{}lower\_\allowbreak{}limit} & (Background-subtracted mosaic only) The inner edge of the background region, as a signed $\Delta r$ in km; rows between $\Delta r$ = −1000 km and this value are used for background modeling \\
\hline
\code{bkgnd\_\allowbreak{}upper\_\allowbreak{}limit} & (Background-subtracted mosaic only) The outer edge of the background region, as a signed $\Delta r$ in km; rows between this value and $\Delta r$ = +1000 km are used for background modeling \\
\hline
\end{longtable}
\endgroup
